\chapter*{Resumen}
\addcontentsline{toc}{chapter}{Resumen}

\azul{El resumen constituye una síntesis estructurada del trabajo de titulación. Este apartado aborda, en una extensión sugerida de 250 a 300 palabras, el contexto del problema, el objetivo principal, la metodología ingenieril aplicada y los resultados o conclusiones más relevantes. La literatura metodológica sostiene que un resumen riguroso permite a otros investigadores evaluar rápidamente la pertinencia del estudio \cite[p. ~312]{hernandez2014metodologia}. Para mantener la formalidad, se sugiere redactar el texto en un único párrafo continuo, omitiendo el uso de citas bibliográficas, ecuaciones complejas o referencias a figuras internas.}

\moradoc{Ejemplo: El presente proyecto de título postula el diseño e implementación de un sistema telemétrico para la optimización de recursos hídricos en comités de agua rural de la provincia de Osorno. Históricamente, la administración de estos recintos carece de monitoreo continuo, lo que podría indicar ineficiencias operativas y pérdidas no contabilizadas. Para mitigar esta problemática, la investigación aplicó el marco de trabajo CDIO (Concebir, Diseñar, Implementar, Operar), articulando el ciclo de vida del artefacto desde su evaluación en terreno hasta su despliegue operativo. A nivel técnico, el desarrollo integró sensores de caudal basados en hardware de código abierto y una plataforma centralizada bajo una arquitectura de microservicios. Las métricas extraídas durante la fase de pruebas sugieren que el sistema reduce el tiempo de respuesta ante fugas estructurales en un 40\%, optimizando el control del caudal. Finalmente, la solución tecnológica implementada discrepa de los procedimientos manuales tradicionales, proporcionando una herramienta escalable que fortalece la toma de decisiones estratégicas en las comunidades rurales.}

\vspace{1cm}

\textbf{Palabras clave:} \azul{Se requieren entre tres y cinco términos técnicos, ordenados alfabéticamente y separados por comas, que encapsulen los dominios principales del proyecto. Estos descriptores facilitan la futura indexación del documento en repositorios académicos.}

\moradoc{Ejemplo: Arquitectura de microservicios, CDIO, Gestión hídrica, Sistemas distribuidos, Telemetría.}